\documentclass[12pt,a4paper]{article}
\usepackage[T1]{fontenc}
\usepackage[utf8]{inputenc}
\usepackage[margin=25.4mm]{geometry}
\usepackage{graphicx,booktabs,amsmath,amssymb,microtype}
\usepackage[hidelinks]{hyperref}
\usepackage{caption,titlesec}
\titleformat{\section}{\normalsize\bfseries}{}{0pt}{}
\titlespacing*{\section}{0pt}{8pt}{5pt}
\captionsetup{font=normalsize,labelfont=bf,skip=4pt,hypcap=false}
\setlength{\parindent}{0pt}
\setlength{\parskip}{5pt}
\setlength{\emergencystretch}{2em}
\graphicspath{{assets/}}
\newcommand{\fig}[3]{\par\begingroup\centering\includegraphics[width=#1\linewidth]{#2}\captionof{figure}{#3}\par\endgroup\vspace{3pt}}
\newcommand{\sub}[1]{\par\textbf{#1}\par}
\begin{document}
\begin{center}\textbf{Pool-ball localisation and recognition}\\
Manon Yiqin Galian Ng (23375206)\end{center}

\section*{1. Finding the balls}
The system finds possible balls, recognises them and converts their centres to
table coordinates. AllBalls supplies training examples; Table1--Table6 contain
90 annotated balls. One parameter set processes every image. Later rules followed
inspection of these six images, so the scores are not an independent test on
new scenes.

\fig{.92}{pipeline.pdf}{The processing stages. Dashed arrows supply AllBalls models. Position conversion is separate from recognition.}

\sub{Masking and circle detection}
The camera does not move, so the supplied table corners define a fixed playing
area: $(337,35)$, $(1570,21)$, $(1888,451)$ and $(38,472)$, written as
$(\text{column},\text{row})$. Pixels outside this quadrilateral are masked out.
This prevents the surrounding room from producing ball detections. A
$5\times5$ Gaussian filter, with standard deviation 1.5 pixels, reduces noise
in the greyscale image before circular Hough detection.

Hough settings were compared on AllBalls in a 480-combination search. The
selected accumulator resolution ratio is 1.5, minimum centre spacing 18 pixels,
upper edge threshold 75 and circle-voting threshold 35; allowed radii are
15--30 pixels. The circle estimates are candidates: pockets and parts of
neighbouring balls can also produce circular edges.

\sub{Rejecting pockets}
Candidate centres are converted to table coordinates and compared with six
known pocket positions. A candidate within 25 table units of a pocket centre
is rejected. This uses the table layout to remove predictable false detections,
but could also remove a genuine ball close to a pocket.

\fig{.83}{detection_stages.pdf}{AllBalls detection stages. Cyan guides mark centre rows used to estimate radius. Recognition uses original colours.}

\newpage
\section*{2. Recognising ball type and colour}
Each detected circle defines a square region of interest (ROI). A circular mask
restricts recognition to pixels inside the estimated ball. The system asks two
separate questions: \emph{how much of the ball appears white?} and \emph{which
body colour does it resemble?} This separates type (cue, solid or striped) from
colour, since solid and striped partners share a colour.

\sub{K-means estimates the amount of white}
K-means divides the ROI's BGR pixels into two colour groups. Pixels in a group
need not touch each other. Each group is represented by its average colour,
which is converted to HSV. Saturation $S$ describes colour strength; value $V$
describes brightness. Both use the range 0--255 here.

A group is light if $V\geq170$ and $S\leq165$. Otherwise it is dark if
$V\leq90$, and coloured otherwise. The sizes of these groups give light and
non-light percentages. At least 85\% light initially suggests cue; at least
70\% non-light suggests solid; a remaining mixture suggests stripe. Dark pixels
count as non-light so dark ball bodies are not discarded. Clustering uses three
initialisations, a fixed seed, and at most 50 iterations or a centre change
below 0.2. It reduces sensitivity to individual pixels, but may merge a pale
body colour with white.

\fig{.91}{recognition_stages.pdf}{Solid 1 (top) and striped 9 (bottom). The light-cluster mask estimates white coverage; the body mask supplies colour samples.}

\sub{HSV histograms identify the body colour}
AllBalls supplies eight colour models: yellow, blue, red, purple, orange, green,
maroon and black. Solid/striped partners contribute equally to their shared
model; cue is excluded. Before measuring body colour, the circular mask is
eroded with a $3\times3$ ellipse and light pixels are removed. This reduces
white-marking and background influence. Fewer than 12 remaining pixels gives
no usable body-colour evidence.

Each histogram counts pixels in $18\times8\times8$ hue/saturation/value bins,
then divides by its total. The bin ranges are $[0,180)$, $[0,256)$ and $[0,256)$.
Neighbouring bins are smoothed with weights $(1,2,1)/4$; hue wraps around.
For nearly grey pixels (BGR channel range below 15), hue is set to zero.
A crop is compared with each trained histogram using one minus Bhattacharyya
distance: a higher score means a closer colour distribution~\cite{colour}.
Backprojection shows where a model responds in the crop; its brightness is
scaled for display and is not the score used to choose an identity.

\newpage
\section*{3. Combining evidence and producing positions}
\sub{Avoiding conflicting identities}
The initial type is useful but can be wrong near a threshold. For example, a
mostly white stripe may be labelled cue. We therefore keep scores for possible
identities instead of immediately accepting one label. A numbered-ball score
multiplies colour similarity by agreement with the expected white coverage:
cue 85--100\%, solid 0--30\%, stripe 30--85\%. Agreement is highest inside its
range and falls smoothly outside it, using a Gaussian width of 0.10 in light
fraction. Cue's score uses light fraction in place of histogram similarity;
black supports ball 8 only. These scores are not calibrated probabilities.

The Hungarian method~\cite{assignment} maximises total score while using each
identity at most once. If two candidates prefer cue, one may receive its
next-best striped identity. Balls need not all be present. Unknown has effectively
zero score and is available when evidence is missing or uniqueness leaves no
supported choice. Receiving an identity requires usable body colour or an
initial cue label. Unknown retains the detection, including a false circle.

\sub{Padding white candidates and selecting the cue}
Hough sometimes follows a stripe's inner white region rather than its full
boundary. Every ROI initially labelled cue is therefore expanded by one pixel,
and its type and histogram evidence are recomputed. This may expose a coloured
rim, but also admits cloth. Other ROIs are unchanged.

The full cue should generally be larger than an internal white patch, after
accounting for depth. A median-slope line~\cite{sen} fits the 16 AllBalls centre
rows and radii (Appendix B):
\[
 \hat r(y)=\max(1,20.274+0.009677y).
\]
With at least two original cue-like candidates, cue is reserved for the largest
original-radius-to-expected-radius ratio, $r_i/\hat r(y_i)$. Padding does not
increase this ratio. This is a forced rule, not proof that the chosen ball is cue.

\fig{.82}{row_radius_fit.pdf}{The radius model uses a narrow training band. Dashed sections estimate sizes beyond the observed rows. Source measurements appear in Appendix B.}

\sub{Converting centres to table coordinates}
A homography maps the supplied corners onto a $1434\times683$ image, matching
the table dimensions in millimetres. The same perspective correction converts
detected centres. Pixel endpoints are 0--1433 and 0--682. Recognition leaves
positions unchanged. Appendix A shows the final views.

\newpage
\section*{4. Measuring performance}
Predicted centres are matched one-to-one with annotated centres within 15 table
units, without using ball identities. The matching finds as many valid pairs
as possible, considering nearby alternatives first. A matched detection is a
true positive (TP); an unmatched detection is a false positive (FP); an unmatched
annotation is a false negative (FN). The matching does not guarantee the lowest
possible total position error in a crowded group.
\[
 P=\frac{TP}{TP+FP},\qquad R=\frac{TP}{TP+FN},\qquad F1=\frac{2PR}{P+R}.
\]
Precision measures how many detections are real balls; recall measures how many
annotated balls were found. Mean position error averages distances over matched
pairs. Identity accuracy counts correct identities among those matches. Unknown
counts as an error on a genuine ball. End-to-end identification also includes
missed balls, using all annotations as its denominator.

\begin{center}
\captionof{table}{Results for each image. Error is in table units (approximately mm); identity accuracy considers matched detections only.}
\begin{tabular}{lrrrrrrr}
\toprule
Image & GT & TP & FP & FN & Error & Correct ID & ID acc.\\
\midrule
Table1 & 16 & 14 & 1 & 2 & 7.68 & 11/14 & 78.57\% \\
Table2 & 16 & 16 & 1 & 0 & 3.58 & 16/16 & 100.00\% \\
Table3 & 16 & 15 & 0 & 1 & 3.12 & 15/15 & 100.00\% \\
Table4 & 15 & 14 & 0 & 1 & 3.89 & 14/14 & 100.00\% \\
Table5 & 14 & 14 & 0 & 0 & 5.83 & 13/14 & 92.86\% \\
Table6 & 13 & 13 & 0 & 0 & 5.76 & 11/13 & 84.62\% \\
ALL & 90 & 86 & 2 & 4 & 4.91 & 80/86 & 93.02\% \\
\bottomrule
\end{tabular}

\end{center}

Detection precision is $86/88=97.73\%$, recall $86/90=95.56\%$ and F1
$96.63\%$. Mean error is 4.91 units, averaged over all 86 matches. Type and
identity totals both equal 80/86 (93.02\%). End-to-end identification is
$80/90=88.89\%$. The two Unknowns are Table1's real stripe 14 and Table2's
unmatched extra circle. AllBalls obtains 16/16 identities, but it supplied the
training examples, so this is not an independent test result.

\begin{center}
\captionof{table}{Adding recognition stages while keeping the same detections.}
\begin{tabular}{lr}
\toprule Configuration & Correct identities\\\midrule
K=2 with fixed colour rules & 68/86\\
K=2 with HSV histograms & 72/86\\
Add unique assignment & 75/86\\
Add selective +1 pixel padding & 79/86\\
Reserve largest depth-adjusted cue & 80/86\\
\bottomrule\end{tabular}
\end{center}

These gains change recognition, not detection counts. The same settings were
used throughout each run, and annotations were loaded after predictions were
made. However, choosing rules after inspecting Table1--Table6 still influences
the method. New images would be needed to estimate performance independently.

\newpage
\section*{5. Explaining the results}
\fig{.92}{recognition_cases.pdf}{Examples of corrections and remaining errors. Arrows show earlier-to-selected identities; circles show original detections.\label{fig:cases}}

Padding corrects cue/stripe conflicts in Tables3 and 6. Extra colour lets the
striped candidate take another identity, releasing cue for the real cue ball.
Two corrected pairs give four correct identities. The crop alone may still
prefer cue; the joint assignment resolves the conflict.

Table1's real cue includes misleading green cloth. Stripe 14 has only eight
usable coloured pixels after padding, below the required 12. Reserving cue
restores the real cue but leaves stripe 14 Unknown. Table1 also swaps solid 5
and stripe 9.

Tables5 and 6 confuse blue stripes with solid 4; Table6 calls solid 3 stripe 11.
Orientation and highlights change white coverage; inaccurate circles include
cloth or neighbours. Uniqueness can also create wrong swaps or give a false
circle a plausible identity (Figure~\ref{fig:cases}(f)).

\fig{.70}{detection_failures.pdf}{Yellow circles are detections; red crosses are supplied centres projected into the source image.}

Four balls are missed: 6 and 7 (Table1), and 10 (Tables3/4). Occlusion and merged
boundaries weaken circular evidence; minimum spacing may suppress close centres.

\newpage
\section*{6. Alternatives and limitations}
\sub{Why additional segmentation was not retained}
Watershed separates regions by growing them from starting markers~\cite{watershed}.
It was tested to separate overlapping balls. Starting from Hough markers did
not recover any misses: a missing ball often also lacked a useful marker.
Finding clusters first and applying more permissive local Hough searches
recovered three annotations, but lost three previous matches and increased
false positives from 2 to 16. F1 fell from 96.63\% to 89.58\%.

Connected-component analysis (CCA) measured connected white/coloured areas and
how much of the crop boundary they reached. The aim was to distinguish a stripe
from a small white number patch. With the same HSV colour models, the tested
rules reduced correct identities from 72 to 69. White-region shape changes with
orientation, and highlights or shadows can split one region into several.

Chamfer matching compares a proposed circle boundary with nearby image
edges~\cite{chamfer}. Searching centre shifts up to two pixels and radius growth
up to eight fixed the Table1 cue/14 swap when an outer supported circle was
preferred. However, it scored 77/86 overall, below padding's 79/86. A one-pixel
shift exposed 15 body pixels instead of eight: a useful correction, but sensitive
to crop placement and the 12-pixel cutoff. Direct HSV pixel counting also scored
79/86 in the final combination, compared with K-means' 80/86.

\fig{.91}{alternatives.pdf}{Watershed and chamfer examples. A plausible region boundary does not necessarily improve the final detection or identity.}

\sub{Ground truth and generalisation}
Annotations supply centres and identities, not visible masks or true image
radii. They cannot establish whether a segmentation boundary is correct.
Occlusion and shadows make some centres difficult to judge; no specific
annotation has been shown incorrect. The homography describes the cloth plane,
whereas balls extend above it. The 15-unit tolerance is about 35\% of the 43 mm
ball diameter, so checking other tolerances would help assess the results.

The fixed camera, narrow radius range and limited lighting variation favour
this system. AllBalls covers only rows 143--197, making much of the radius
model an extrapolation. Cue reservation can fail if a striped patch is larger
or the real cue was not initially labelled white. Repeated positions and a
shared viewpoint also limit the diversity of the six images. More varied
training views and new test scenes are needed. The observed 80/86 is promising,
but does not demonstrate the same accuracy under different conditions.

\newpage
\section*{Appendix A: final results for all six tables}
\label{app:results}
These are the final rectified views from the selected configuration, using the
same parameters for every image. Labels show predictions, including Unknown;
they have not been replaced with the annotated identities. Positions use the
original Hough centres. Individual full-resolution images accompany the report
source; the overview below retains all six images for comparison.

\fig{1}{final_tables.pdf}{Final predicted identities and positions for Table1--Table6. Read together with the per-image metrics and failure examples in the main text.}

All six real cue balls are correctly identified. Table1 contains three identity
errors among its 14 matched detections; Table5 contains one, and Table6 two.
Tables2--4 correctly identify every matched detection, although Tables3 and 4
each miss one ball. The extra Table2 detection remains Unknown and still counts
as a false positive.

\newpage
\section*{Appendix B: radius measurements and resources}
\fig{.92}{row_radius_samples.pdf}{The AllBalls centre-row band and enlarged measured circles. Crosses mark centres; numbers identify detections rather than ball identities. Circles supply the row/radius pairs used in the model.}

The slope is the median of 71 pairwise slopes with centre rows at least ten
pixels apart. The intercept is the median of the 16 values $r_i-by_i$. Medians
reduce the influence of unusually large or small detections, but cannot remove
a shared detection bias or create evidence at unobserved depths.

\sub{Resources and generative AI acknowledgement}
Images, dimensions, corners and annotations are supplied module materials.
OpenAI ChatGPT/Codex assisted with design, C++ experiment/evaluation code,
debugging, comparisons, figure scripts, diagrams and this report's prose.
AI-generated contributions include histogram, identity, padding, cue-size and
alternative-method experiments, numerical summaries and draft illustrations.

\begin{thebibliography}{5}
\raggedright
\setlength{\itemsep}{0pt}\setlength{\parskip}{0pt}
\bibitem{colour} M. Swain and D. Ballard. Color indexing. \emph{IJCV} 7, 11--32 (1991).
\href{https://doi.org/10.1007/BF00130487}{doi:10.1007/BF00130487}.
\bibitem{assignment} H. Kuhn. The Hungarian method for the assignment problem.
\emph{Naval Research Logistics Quarterly} 2, 83--97 (1955).
\bibitem{sen} P. Sen. Estimates of the regression coefficient based on Kendall's tau.
\emph{JASA} 63, 1379--1389 (1968).
\bibitem{watershed} L. Vincent and P. Soille. Watersheds in digital spaces.
\emph{IEEE TPAMI} 13, 583--598 (1991).
\bibitem{chamfer} G. Borgefors. Hierarchical chamfer matching: a parametric edge
matching algorithm. \emph{IEEE TPAMI} 10, 849--865 (1988).
\end{thebibliography}
\end{document}
